%%=============================================================================
%% Methodologie
%%=============================================================================

\chapter{\IfLanguageName{dutch}{Methodologie}{Methodology}}%
\label{ch:methodologie}

%% TODO: In dit hoofstuk geef je een korte toelichting over hoe je te werk bent
%% gegaan. Verdeel je onderzoek in grote fasen, en licht in elke fase toe wat
%% de doelstelling was, welke deliverables daar uit gekomen zijn, en welke
%% onderzoeksmethoden je daarbij toegepast hebt. Verantwoord waarom je
%% op deze manier te werk gegaan bent.
%% 
%% Voorbeelden van zulke fasen zijn: literatuurstudie, opstellen van een
%% requirements-analyse, opstellen long-list (bij vergelijkende studie),
%% selectie van geschikte tools (bij vergelijkende studie, "short-list"),
%% opzetten testopstelling/PoC, uitvoeren testen en verzamelen
%% van resultaten, analyse van resultaten, ...
%%
%% !!!!! LET OP !!!!!
%%
%% Het is uitdrukkelijk NIET de bedoeling dat je het grootste deel van de corpus
%% van je graduaatsproef in dit hoofstuk verwerkt! Dit hoofdstuk is eerder een
%% kort overzicht van je plan van aanpak.
%%
%% Maak voor elke fase (behalve het literatuuronderzoek) een NIEUW HOOFDSTUK aan
%% en geef het een gepaste titel.
Dit hoofdstuk beschrijft hoe Forcebit Ticketing werd aangepakt. Het oorspronkelijke plan van aanpak was een compacte planning voor een eerste prototype. Tijdens de uitvoering werd die planning aangepast wanneer mentorfeedback, technische keuzes of onderhoudbaarheid daar aanleiding toe gaven. Per fase worden doel, werkwijze en resultaat samengevat.

\section{Oorspronkelijke planning}

Het oorspronkelijke plan vertrok van een korte ontwikkelperiode waarin eerst de basisstructuur en daarna de belangrijkste functionaliteit werd gebouwd:

\begin{center}
\footnotesize
\begin{tabular}{@{}p{0.47\linewidth}p{0.47\linewidth}@{}}
\textbf{Dag 1:} requirements, databankontwerp en wireframes & \textbf{Dag 6:} React Native login, registratie en navigatie \\
\textbf{Dag 2:} backendproject, MySQL databank en modellen & \textbf{Dag 7:} klantenschermen voor overzicht, aanmaken en detail \\
\textbf{Dag 3:} authenticatie en rollen & \textbf{Dag 8:} adminschermen \\
\textbf{Dag 4:} ticket CRUD & \textbf{Dag 9:} bestandsuploads en e-mailnotificaties \\
\textbf{Dag 5:} berichten en antwoorden & \textbf{Dag 10:} testen, fouten oplossen, documentatie en demo
\end{tabular}
\end{center}

De uitvoering volgde deze planning op hoofdlijnen, maar sommige onderdelen kregen meer aandacht dan eerst voorzien. Vooral chatweergave, zoekfilters, paginatie, responsive layout en documentatie werden later uitgebreid.

\Needspace{2\baselineskip}
\section{Fase 1: scope en analyse}

\textbf{Doel.} De eerste fase moest duidelijk maken welke rollen, gegevens en workflows nodig waren.\par\noindent\textbf{Werkwijze.} Hiervoor werden requirements en wireframes opgesteld vanuit drie perspectieven: klant, administrator en systeem.\par\noindent\textbf{Resultaat.} De kern werd vastgelegd als gebruikers, tickets, berichten en bijlagen. De statusflow werd later verfijnd van \texttt{Open}, \texttt{In Progress} en \texttt{Closed} naar \texttt{New}, \texttt{Open} en \texttt{Closed}, waarbij \texttt{New} alleen voor pas aangemaakte tickets dient.

\Needspace{2\baselineskip}
\section{Fase 2: technische basis}

\textbf{Doel.} De tweede fase moest een backend opleveren waarop de ticketflow betrouwbaar kon steunen.\par\noindent\textbf{Werkwijze.} Er werd een prototype gebouwd met ASP.NET Core, Entity Framework Core en MySQL. De code werd meteen gelaagd opgezet om verantwoordelijkheden te scheiden.\par\noindent\textbf{Resultaat.} De belangrijkste deliverables waren entities, DTOs, controllers, repositories, services, configuratie, migraties, JWT-authenticatie, centrale foutafhandeling en het Options pattern \autocite{MicrosoftOptionsPattern2026}.

\Needspace{2\baselineskip}
\section{Fase 3: frontend en feedback}

\textbf{Doel.} De derde fase moest de backend bruikbaar maken voor klanten en administrators.\par\noindent\textbf{Werkwijze.} De wireframes werden omgezet naar frontend-schermen. Expo werd gebruikt als ontwikkelplatform \autocite{ExpoDocs2026}. Formik werd gebruikt voor formulierstate en Yup voor schema-validatie \autocite{FormikDocs2026,YupDocs2026}. Herbruikbare logica werd verplaatst naar hooks, API modules, componenten, stijlen en utilities.\par\noindent\textbf{Resultaat.} De basisflow werd aangevuld met mentorfeedback: chatbubbels, afbeeldingspreview, uploadlimiet, responsive layout, browsernavigatie, zoekvelden, filters en paginatie.

\Needspace{2\baselineskip}
\section{Fase 4: controle en documentatie}

\textbf{Doel.} De laatste fase moest nagaan of de applicatie ook bij realistische data en langere gesprekken bruikbaar en onderhoudbaar bleef.\par\noindent\textbf{Werkwijze.} Er werd realistische testdata toegevoegd voor meerdere klanten, tickets en langere gesprekken. Daarna werden zoekfilters, statusgroepen, paginatie, database-startup en documentatie gecontroleerd.\par\noindent\textbf{Resultaat.} Code, comments en README-bestanden werden opgeschoond. Zo sluit deze fase aan bij de onderzoeksvraag rond onderhoudbaarheid: de applicatie blijft niet alleen werkend, maar ook uitbreidbaar en onderhoudbaar.